\documentclass[11pt,a4paper]{article}
\usepackage[T1]{fontenc}
\usepackage[utf8]{inputenc}
\usepackage{amsmath,amssymb,amsthm,newpxtext,newpxmath,microtype}
\usepackage[a4paper,margin=28mm]{geometry}
\usepackage{booktabs,array,longtable}
\usepackage[colorlinks=true,linkcolor=blue,citecolor=blue,urlcolor=blue]{hyperref}
\hypersetup{pdftitle={Polynomial relations in totient kernels and separation of dilated residue values},
 pdfsubject={Round-six ordinary proof candidates, not Lean compiled},
 pdfauthor={Research memorandum, 7 September 2026}}
\usepackage{fancyhdr}
\pagestyle{fancy}\fancyhf{}\fancyhead[L]{\small Totient kernels and dilated residue values}\fancyhead[R]{\small\thepage}
\setlength{\headheight}{14pt}\renewcommand{\headrulewidth}{0pt}
\newtheorem{theorem}{Theorem}[section]
\newtheorem{lemma}[theorem]{Lemma}
\newtheorem{corollary}[theorem]{Corollary}
\newtheorem{proposition}[theorem]{Proposition}
\theoremstyle{definition}\newtheorem{definition}[theorem]{Definition}
\theoremstyle{remark}\newtheorem{remark}[theorem]{Remark}
\newcommand{\Q}{\mathbb Q}\newcommand{\R}{\mathbb R}\newcommand{\Z}{\mathbb Z}\newcommand{\N}{\mathbb N}
\newcommand{\ph}{\varphi}\newcommand{\rad}{\operatorname{rad}}
\newcommand{\Span}{\operatorname{span}}\newcommand{\dist}{\operatorname{dist}}
\newcommand{\ev}{\operatorname{ev}}\newcommand{\ord}{\operatorname{ord}}
\newcommand{\eps}{\varepsilon}
\title{Polynomial relations in totient kernels\\and separation of dilated residue values}
\author{Round-six research memorandum}
\date{7 September 2026}
\begin{document}
\maketitle
\begin{abstract}
The existing finite-level totient basis admits a proposed polynomial normal form.
Over $\Q(n)$, all relations among its $k^e+1$ coordinates are generated by one
homogeneous polynomial of degree $2^{\omega(k)}$, involving only the two zero
channels. At base two this polynomial is $(Y-X)(Y-2X)$.
An elementary local-density argument proves completeness and shows that algebraic
independence of nonproportional affine sections survives every $o(n)$ perturbation.
Consequently the supplied rational control retains that independence.
A separate isolated-pulse construction proves joint rationality classification
for dyadic residue observables at any finite collection of bases $b^d$.
The new arguments are ordinary proofs submitted for independent checking; none
has been compiled in Lean. They do not prove irrationality of the unreduced series.
\end{abstract}

\section{The polynomial normal form}
For $k\ge2$ and $e\ge1$, use the existing canonical coordinates
\[
 X(n)=\ph(n),\qquad Y(n)=\ph(kn),\qquad
 Z_{j,r}(n)=\ph(k^j n+r),
 \quad 1\le j\le e,\ 0<r<k^j,\ k\nmid r.
\]
There are $q=k^e+1$ coordinates. A polynomial with coefficients in $\Q(n)$
is evaluated outside the finitely many poles of its coefficients. Its evaluation
is zero \emph{eventually} when it vanishes for all sufficiently large integers $n$.

Let $\mathcal P(k)$ be the set of prime divisors of $k$. For $T\subseteq\mathcal P(k)$ put
\[
 \lambda_T=k\prod_{p\in\mathcal P(k)\setminus T}(1-1/p)
 =\ph(k)\prod_{p\in T}\frac p{p-1},\qquad h=2^{\omega(k)},
\]
and define
\begin{equation}\label{eq:Pk}
 P_k(X,Y)=\prod_{T\subseteq\mathcal P(k)}(Y-\lambda_T X).
\end{equation}
Every $\lambda_T$ is a positive integer.

\begin{theorem}[Complete polynomial normal form]\label{thm:kernel-algebra}
The ideal of eventual polynomial relations over $\Q(n)$ among
$X,Y,(Z_{j,r})$ is the principal ideal $(P_k)$.
Equivalently, every polynomial expression has a unique normal form of degree
less than $h$ in $Y$, and that normal form evaluates to zero eventually only
when all its coefficients vanish.
For the full kernel truncation, first apply the scalar reductions in
\texttt{thm:kkernelrank}; the only remaining polynomial relation is \eqref{eq:Pk}.
\end{theorem}

The proof is in Section~\ref{sec:complete}. Its additional arithmetic ingredient
is an open set of limiting values of normalised nonproportional affine totients,
including on any fixed progression. This prevents hidden polynomial relations
on each local divisibility class.

At base two, the complete relation is
\begin{equation}\label{eq:dyadicpoly}
 (Y-X)(Y-2X)=0,\qquad Y^2=3XY-2X^2.
\end{equation}
At base six it is
\[
 (Y-2X)(Y-3X)(Y-4X)(Y-6X)=0.
\]
The two zero coordinates are linearly independent despite satisfying a quadratic
relation in the dyadic case. Their ratio is selected by the parity of $n$.

\begin{corollary}[Ranks of products]\label{cor:hilbert}
Let $W_D$ be the rational span of all products of total degree at most $D$ in
the $q=k^e+1$ canonical coordinates, including the empty product $1$.
Then
\[
 \dim_\Q W_D=\binom{q+D}{q}-\binom{q+D-h}{q},
 \qquad h=2^{\omega(k)},
\]
where the second binomial coefficient is zero when $D<h$.
The homogeneous Hilbert series is $(1-t^h)/(1-t)^q$.
The same normal monomials are independent over $\Q(n)$.
\end{corollary}
\begin{proof}
The polynomial $P_k$ is homogeneous of degree $h$ and monic in $Y$.
Polynomial division reduces uniquely to monomials with $Y$-exponent less than $h$.
Theorem~\ref{thm:kernel-algebra} proves their independence, including over $\Q(n)$.
Counting monomials of degree at most $D$, and subtracting those divisible by
$Y^h$, gives the formula. Monic division takes place over $\Q$ when the original
coefficients are constant, so it also proves the asserted $\Q$-dimension.
\end{proof}

For $(k,e)=(2,3)$, quadratic products and constants have rank $54$ among $55$
monomials. For $(k,e)=(6,1)$, all $120$ monomials of degree at most three are
independent; the $330$ monomials through degree four have rank $329$.
These predictions are independently checked by exact finite-field evaluation
in the accompanying script, which is evidence for finite instances rather than
a proof of the theorem.

\paragraph{Depth and arithmetic weight.}
The presentation also gives a compatible polynomial extension
\[
 \mathcal A_{k,e+1}\simeq
 \mathcal A_{k,e}[U_1,\ldots,U_{k^{e+1}-k^e}],
\]
where $\mathcal A_{k,e}$ is the algebra generated by the truncated kernel over
$\Q(n)$. The branching degree $2^{\omega(k)}$ stays fixed as depth increases.
For $f(n)=n^a\ph(n)^m$, $a\ge0$, $m\ge1$, the same proof gives the analogous
presentation, with roots $k^a\lambda_T^m$ in place of $\lambda_T$.
These roots are distinct because the $\lambda_T$ are distinct and positive.
Algebraic independence of the nonproportional $f$-sections follows by the
injective substitution $V_i\mapsto L_i(n)^a V_i^m$ in the algebraically
independent totient sections. Their scalar reductions are $k^{at}C_k(t,u)^m$.
The hypothesis $m\ge1$ matters: pure powers have polynomial sections in $n$.

\section{An elementary open-box construction}\label{sec:box}
The following form of density is sufficient; a complete description of the
limiting support is unnecessary.

\begin{lemma}[Positive-density approximation in a box]\label{lem:box}
Let $L_i(n)=a_i n+b_i$ ($1\le i\le s$), where $a_i\ge1$, $b_i\ge0$, and
$a_i b_j-a_j b_i\ne0$ for $i\ne j$.
There are $c_i>0$ such that, for every $x_i\in(0,c_i)$ and every $\eps>0$,
\[
 \left\{n:\max_i\left|\frac{\ph(L_i(n))}{L_i(n)}-x_i\right|<\eps\right\}
\]
has positive lower natural density. Initial zero arguments are omitted.
The same statement holds with $n$ restricted to any prescribed nonconstant
arithmetic progression, with constants that may depend on that progression.
\end{lemma}
\begin{proof}
Choose $Y_0>s$ larger than every prime dividing a slope or a nonzero cross
determinant. Fix a residue $a_0$ modulo the product of primes at most $Y_0$;
for example $a_0=1$ is permissible. Put
\[
 c_i=\prod_{\substack{p\le Y_0\\p\mid L_i(a_0)}}(1-1/p)>0.
\]
These constants are fixed before the accuracy is chosen.
Fix $x_i\in(0,c_i)$ and a desired error $\eps$.
Choose $\eta>0$ small enough that $2\eta<\eps$, and then choose $Y>Y_0$
so large that $s/(Y-1)<\eta/4$ and $1/Y<\eta$.

For each $i$, choose a finite set $\mathcal Q_i$ of primes greater than $Y$,
the sets mutually disjoint, such that
\[
 \left|c_i\prod_{p\in\mathcal Q_i}(1-1/p)-x_i\right|<\eta.
\]
To justify this without a density theorem, take unused primes in increasing
order until their Euler product first crosses $x_i/c_i$.
The product tends to zero because $\sum_p1/p$ diverges; its final relative
step is at most $1/Y$. Thus the overshoot in absolute value is less than
$1/Y$, and $Y$ may be enlarged initially to meet the displayed accuracy.
Removing the finite sets already used does not affect this argument.

Prescribe $n\equiv a_0$ at primes at most $Y_0$.
At every prime $Y_0<p\le Y$, choose $n$ so that none of the $L_i(n)$ is
zero modulo $p$. At most $s<p$ classes are forbidden.
For $p\in\mathcal Q_i$, prescribe $L_i(n)\equiv0\pmod p$.
All slopes are units modulo these primes. The distinct-root property follows
from the nonzero cross determinants, so this prescription prevents $p$ from
dividing every other $L_j(n)$.
CRT gives one progression $n=a+Qt$ on which all prescribed factors hold.
Let
\[
 C_i=c_i\prod_{p\in\mathcal Q_i}(1-1/p).
\]
Every prime dividing $Q$ has now been accounted for in $C_i$. Consequently
\[
 \frac{\ph(L_i(a+Qt))}{L_i(a+Qt)}
 =C_i\prod_{\substack{p>Y,\ p\nmid Q\\p\mid L_i(a+Qt)}}(1-1/p).
\]
For $1\le t\le T$, the affine arguments are at most $C(T+1)$ for a fixed
constant $C$ depending on this progression. If $p\nmid Q$ and $p>Y$, its
slope is nonzero modulo $p$, so there are at most $T/p+1$ relevant values of $t$.
It follows that
\begin{align*}
 \frac1T\sum_{t=1}^T\sum_{i=1}^s
 \sum_{\substack{p>Y,p\nmid Q\\p\mid L_i(a+Qt)}}\frac1p
 &\le s\sum_{p>Y}\frac1{p^2}
   +\frac{s}{T}\sum_{p\le C(T+1)}\frac1p\\
 &\le \frac{s}{Y-1}+O_{a,Q,L_i}\!\left(\frac{\log(T+2)}T\right).
\end{align*}
For all sufficiently large $T$, this is below $\eta/2$.
At least half the parameters $t$ therefore have their combined tail sum
below $\eta$. For those parameters, the elementary inequality
$\prod(1-u_j)\ge1-\sum u_j$ for $0\le u_j\le1$ gives an error at most
$\eta$ in every coordinate. Combining this with $|C_i-x_i|<\eta$ gives
the required approximation on a set of lower density at least $1/(2Q)$.

For a prescribed progression $n=a+Mt$, first shift $t$ enough to make all
intercepts nonnegative, then apply the same argument to $L_i(a+Mt)$.
Their cross determinants are $M(a_i b_j-a_j b_i)\ne0$.
\end{proof}

\begin{remark}
The proof imposes no primality condition on any affine value. The arbitrarily
large primes are local divisors. Uncontrolled Euler factors are bounded in mean.
The constants $c_i$ do not shrink when the requested error decreases: primes
between $Y_0$ and $Y$ are avoided rather than incorporated into $c_i$.
This order of choices is necessary for obtaining one fixed open box.
\end{remark}

\begin{theorem}[Algebraic independence, stable under $o(n)$]\label{thm:stable}
Under the hypotheses of Lemma~\ref{lem:box}, suppose that real sequences $G_i(n)$
satisfy $G_i(n)-\ph(L_i(n))=o(n)$. Then no nonzero polynomial
\[
 P(U,V_1,\ldots,V_s)\in\R[U,V_1,\ldots,V_s]
\]
can satisfy $P(n,G_1(n),\ldots,G_s(n))=0$ eventually.
In fact, the set on which this evaluation is nonzero has positive lower density.
In particular, the affine totient sequences are algebraically independent over
$\Q(n)$, including after restriction to any fixed progression.
\end{theorem}
\begin{proof}
Let $P_D$ be the highest total-degree homogeneous part of $P$. Define
\[ H(x)=P_D(1,a_1x_1,\ldots,a_sx_s). \]
This polynomial is nonzero because dehomogenisation at $U=1$ is injective
on homogeneous polynomials of a fixed degree.
Choose $x$ in the open box of Lemma~\ref{lem:box} where its value is nonzero,
and a smaller open neighbourhood on which its absolute value is bounded below.
On a positive-lower-density set of sufficiently large $n$, the ratios
$\ph(L_i(n))/L_i(n)$ lie in this neighbourhood.
Since $L_i(n)/n\to a_i$ and $(G_i(n)-\ph(L_i(n)))/n\to0$,
the highest homogeneous contribution to
\[ n^{-D}P(n,G_1(n),\ldots,G_s(n)) \]
is bounded away from zero there. The contributions of lower degree tend uniformly
to zero, because all $G_i(n)/n$ remain bounded. The evaluation is therefore nonzero.
Clearing denominators proves the statement over $\Q(n)$.
For a fixed progression use the last assertion of Lemma~\ref{lem:box}.
\end{proof}

\section{Completeness of the kernel relation}\label{sec:complete}
\begin{proof}[Proof of Theorem~\ref{thm:kernel-algebra}]
First, the affine forms underlying $X$ and the retained $Z_{j,r}$ are pairwise
nonproportional. Distinct positive-residue forms could be proportional only if
$r/k^j=r'/k^{j'}$. At distinct levels this makes the residue at the larger level
divisible by $k$, contrary to retention. At the same level the residues coincide.
No positive-residue form is proportional to $n$.
Thus these $k^e$ sequences are algebraically independent over $\Q(n)$ on every
fixed progression, by Theorem~\ref{thm:stable}.

For $n\ge1$, the local formula for the totient gives
\[
 \frac{Y(n)}{X(n)}
 =k\prod_{\substack{p\mid k\\p\nmid n}}(1-1/p)
 =\lambda_{T(n)},\qquad T(n)=\{p\mid k:p\mid n\}.
\]
Hence $P_k(X(n),Y(n))=0$.
Every subset $T$ occurs on a full progression modulo $\rad(k)$, by prescribing
$n\equiv0\pmod p$ for $p\in T$ and $n\equiv1\pmod p$ otherwise.

The numbers $\lambda_T$ are distinct. Otherwise, cancellation between
$\lambda_T/\lambda_{T'}$ would give equality between two products of
$p/(p-1)$ indexed by disjoint nonempty sets. The largest prime in their union
occurs in one numerator and in no denominator or other numerator, an impossibility.

Now suppose $P(n,X(n),Y(n),Z(n))=0$ eventually, with coefficients in $\Q(n)$.
Restrict to a progression realising $T$. Substitution of $Y=\lambda_T X$
gives a polynomial relation among $X,Z$, so algebraic independence proves
\[
 P(n,X,\lambda_T X,Z)=0
\]
as a polynomial identity. Equivalently $Y-\lambda_T X$ divides $P$.
The distinct linear factors are pairwise nonassociate prime elements in the
polynomial ring over $\Q(n)$. Their product therefore divides $P$.
The reverse implication was proved by the displayed local identity.
Monic division in $Y$ now gives the unique normal form.

Finally, every omitted kernel coordinate is an exact constant multiple of one
retained coordinate, by the live basis theorem. Quotienting the polynomial ring
on all coordinates by these linear reductions identifies it with the polynomial
ring on $X,Y,Z$. Thus the linear reductions and $P_k$ generate the full relation ideal.
\end{proof}

\subsection{The analogous answer for an arbitrary finite affine family}
This gives a useful finite algorithm beyond kernel sections.
Write $a_i n+b_i=c_i L_{j(i)}(n)$ with $L_j$ primitive and distinct.
Put $Q_0=\prod_{p\mid\prod_i c_i}p$, with the empty product equal to one, and
\[
 \lambda_i(a)=\ph(c_i)
 \prod_{\substack{p\mid c_i\\p\mid L_{j(i)}(a)}}\frac p{p-1}
 \quad (a\bmod Q_0).
\]
On the progression $n\equiv a\pmod{Q_0}$,
$\ph(a_i n+b_i)=\lambda_i(a)\ph(L_{j(i)}(n))$.
Define a homomorphism over $\Q(n)$ by
\[
 \psi_a:\Q(n)[Y_i]\longrightarrow\Q(n)[Z_j],
 \qquad Y_i\longmapsto\lambda_i(a)Z_{j(i)}.
\]
The complete ideal of eventual polynomial relations is
\begin{equation}\label{eq:arrangement}
 I=\bigcap_{a\bmod Q_0}\ker\psi_a.
\end{equation}
Indeed, independence on each progression proves necessity, and the local
identities prove sufficiency. Each kernel is generated by linear proportionality
relations within each primitive class. Consequently the closure is a finite
union of rational linear subspaces. Redundant progressions may be merged.
This is an effective intersection description, rather than a claim that one
unreviewed list is a minimal Gr\"obner basis.

For example the four sequences $\ph(n),\ph(2n),\ph(3n),\ph(6n)$ are linearly
independent but satisfy
\[
 \ph(n)\ph(6n)=\ph(2n)\ph(3n).
\]
Their four local lines have direction vectors
$(1,1,2,2),(1,2,2,4),(1,1,3,3),(1,2,3,6)$.
Their ideal is exactly the intersection of the ideals of these four lines.

\section{Value separation at powers of one base}\label{sec:values}
Here the conclusions concern actual real numbers, rather than relations between
coefficient sequences. Fix an integer $b\ge2$ and a finite set $D$ of distinct
positive integers. Define
\[
 T_{m,b^d}(f)=\sum_{n\ge1}f(\ph(n)\bmod m)b^{-dn}.
\]

\begin{theorem}[Joint dyadic rationality classification]\label{thm:joint}
Let $k\ge1$ and let $f_d:\Z/2^k\Z\to\Q$ for $d\in D$. Then
\[
 \sum_{d\in D}T_{2^k,b^d}(f_d)\in\Q
 \quad\Longleftrightarrow\quad
 f_d\text{ is constant on the even residues for every }d\in D.
\]
If these constants are $c_d$, the value is
\[
 \sum_{d\in D}\left[
 f_d(1)(b^{-d}+b^{-2d})+
 c_d\frac{b^{-3d}}{1-b^{-d}}\right].
\]
\end{theorem}

\begin{corollary}[Least-residue values]\label{cor:jointleast}
For every $m\ge3$ and $b\ge2$, the countable family
\[
 \{1\}\ \cup\
 \left\{\sum_{n\ge1}(\ph(n)\bmod m)b^{-dn}:d\ge1\right\}
\]
is linearly independent over $\Q$.
\end{corollary}

\begin{corollary}[All relations between dyadic residue classes]\label{cor:value-dimension}
For $r\bmod2^k$ and $d\in D$, put
\[
 V_{r,d}=\sum_{\ph(n)\equiv r\ (2^k)}b^{-dn}.
\]
The set consisting of $1$ and the $V_{r,d}$ with
$r\in\{2,4,\ldots,2^k-2\}$, $d\in D$, is linearly independent.
It spans every $V_{r,d}$ and $1$, so that span has dimension
\[
 1+|D|(2^{k-1}-1).
\]
\end{corollary}
\begin{proof}
Apply Theorem~\ref{thm:joint} to a proposed relation, taking each observable
to be zero at residue zero and all odd residues. Constancy on the even
residues forces every coefficient to vanish, and then the coefficient of $1$
vanishes. For spanning, the odd-residue-one value is $b^{-d}+b^{-2d}$, every
other odd-residue value is zero, and
\[
 \sum_{r\text{ even}}V_{r,d}=\frac{b^{-3d}}{1-b^{-d}}.
\]
This also treats $k=1$, where the independent family contains only $1$.
\end{proof}

\begin{corollary}[A two-parameter independent family]\label{cor:grid}
Fix $b\ge2$. The numbers $1$ and
\[
 B_{j,d}=\sum_{n\ge1}\operatorname{bit}_j(\ph(n))b^{-dn},
 \qquad j\ge1,\ d\ge1,
\]
are linearly independent over $\Q$. The same holds for $1$ together with
all $A_{2^k,b^d}$ for $k\ge2$, $d\ge1$.
\end{corollary}
\begin{proof}
Any finite relation uses a common dyadic modulus. At each exponent $d$,
Theorem~\ref{thm:joint} makes the corresponding bit combination constant
on even residues. Its value at zero is zero; evaluation at $2^j$ then
recovers and kills each coefficient. Finally
\[
 A_{2^k,b^d}=b^{-d}+b^{-2d}+\sum_{j=1}^{k-1}2^jB_{j,d}
\]
is a triangular change of coordinates at each $d$.
\end{proof}

\subsection{The pulse argument and the smallest dilation}
\begin{lemma}[Signed bounded pulses, any positive integer base]\label{lem:pulse}
Let $a_n\in\Z$, $|a_n|\le C$, and suppose there are arbitrarily long intervals
$[N-L,N+L]$ with $N>L$, $a_N\ne0$ and $a_{N+t}=0$ for $0<|t|\le L$.
Then $\sum_{n\ge1}a_n b^{-n}$ is irrational for every integer $b\ge2$.
\end{lemma}
\begin{proof}
If the value has denominator $q$, every scaled tail
$R_j=\sum_{t\ge1}a_{j+t}b^{-t}$ belongs to $q^{-1}\Z$.
Choose a pulse with $Cb^{-L}/(b-1)<1/q$.
The right zero block gives $|R_N|\le Cb^{-L}/(b-1)<1/q$, hence $R_N=0$.
The left block gives $R_{N-L-1}=a_Nb^{-(L+1)}$, which is nonzero and
has absolute value at most $Cb^{-(L+1)}<1/q$, a contradiction.
\end{proof}

\begin{lemma}[Isolation survives finite dilation mixing]\label{lem:dilation}
Let $m\ge2$. For each $d$ in a nonempty finite set $D\subset\N_{>0}$, let
$g_d:\Z/m\Z\to\Z$ satisfy $g_d(0)=0$. Suppose $d_0=\min D$ and
$g_{d_0}(r)\ne0$ for a residue $r$ with $\gcd(r+1,m)=1$.
Then the integer sequence
\[
 a_N=\sum_{\substack{d\in D\\d\mid N}}
 g_d(\ph(N/d)\bmod m)
\]
has arbitrarily long two-sided isolated nonzero pulses.
\end{lemma}
\begin{proof}
Fix $L$. For every nonzero integer $t$ with $|t|\le L$, choose a distinct
prime $q_t\equiv1\pmod m$ greater than $L$ and $\max D$.
Dirichlet supplies these auxiliary primes.
CRT gives the simultaneous congruences
\[
 p\equiv r+1\pmod m,\qquad
 p\equiv-t d_0^{-1}\pmod{q_t}\quad(0<|t|\le L).
\]
Every residue is a unit: $r+1$ is a unit modulo $m$, and $q_t>|t|,d_0$.
Thus the combined progression is reduced. Dirichlet supplies arbitrarily
large prime solutions $p$, which we require to exceed $\max D$ and to make
$N=d_0p>L+2\max D$.

If $d\in D$ divides $d_0p$, then $\gcd(d,p)=1$ and therefore $d\mid d_0$.
Minimality forces $d=d_0$. Consequently
\[
 a_N=g_{d_0}(\ph(p)\bmod m)=g_{d_0}(r)\ne0.
\]
For a neighbour $N+t$ and any $d\in D$ dividing it, the prime $q_t$ divides
$N+t$ and is coprime to $d$. Hence $q_t\mid(N+t)/d$.
The quotient is positive, and $q_t-1\mid\ph((N+t)/d)$.
As $m\mid q_t-1$, every term in $a_{N+t}$ is $g_d(0)=0$.
This supplies the required pulse.
\end{proof}

\begin{proof}[Proof of Theorem~\ref{thm:joint}]
The displayed value in the constant case follows from $\ph(1)=\ph(2)=1$
and the evenness of $\ph(n)$ for $n\ge3$.
For the converse, subtract that rational constant contribution for each $d$,
and discard the finite terms at $n=1,2$. More explicitly, replace $f_d$ by
an observable equal to $f_d(r)-f_d(0)$ on even residues and zero on odd residues.
This changes the series by a rational number. Clear the common denominators,
then remove the observables that are now identically zero.
If any remain, choose their smallest exponent $d_0$ and an even residue $r$
where its observable is nonzero. The residue $r+1$ is a unit modulo $2^k$.
Absolute convergence allows the finite sum of series to be reindexed as
$\sum_Na_Nb^{-N}$ with the coefficients of Lemma~\ref{lem:dilation}.
They are bounded by the sum of the finite alphabet bounds.
Lemmas~\ref{lem:dilation} and \ref{lem:pulse} make that value irrational,
contradicting rationality of the original combination.
\end{proof}

\begin{proof}[Proof of Corollary~\ref{cor:jointleast}]
Clear denominators in a proposed relation including $1$, and retain only
nonzero coefficients $c_d$. Put $g_d(r)=c_dr$ on the representatives
$0,\ldots,m-1$. At the smallest remaining exponent, choose $r=m-2$.
Then $r\ne0$ and $r+1=m-1$ is a unit modulo $m$.
The two pulse lemmas show that this nonzero combination cannot be rational.
\end{proof}

\subsection{A finite certificate for an individual linear form}
Suppose $a$ has such a pulse at $N$ with half-width $L$, and set
$\eta=Cb^{-L}/(b-1)$. Assume $\eta<1$ and
$b^{-(L+1)}(C+\eta)<1/2$. For any integer $q_0$,
\begin{equation}\label{eq:lower-linear}
 \left|q_0+\sum_{n\ge1}a_nb^{-n}\right|
 \ge b^{-N}(|a_N|-\eta)>0.
\end{equation}
Indeed, multiply by $b^{N-L-1}$. Its integer part is the finite prefix plus
$b^{N-L-1}q_0$. The remaining tail is
$b^{-(L+1)}(a_N+R_N)$, with $|R_N|\le\eta$.
It lies in $(-1/2,1/2)$ and its distance from zero is at least
$b^{-(L+1)}(|a_N|-\eta)$. Distance to the nearest integer is bounded by
absolute value, giving \eqref{eq:lower-linear}.
Thus a proposed integer relation can be refuted by finite prime and congruence
data plus an elementary tail bound, without computing the enormous prefix.
This is not a uniform irrationality measure: the pulse centre depends on the
coefficients and the CRT construction.

\section{What these strengthenings do not prove}\label{sec:boundaries}
\subsection{Algebraic independence of sections is compatible with a rational value}
The live source \texttt{ParityPerturbedRationalControl.lean} constructs a natural
sequence $c$ satisfying
\[
 c(n)=\ph(n)\ (n\text{ odd}),\qquad |c(n)-\ph(n)|\le2,\qquad
 0\le c(n)\le n,\qquad \sum_{n\ge1}c(n)2^{-n}=5/4.
\]
Theorem~\ref{thm:stable} therefore applies to $c(L_i(n))$ for every finite
nonproportional affine family, even on prescribed progressions.
The rational control has joint algebraic independence of those sections over
$\Q(n)$. This strengthens the existing rank-only counterexample.
It does not say that $c$ satisfies the exact scalar reductions of the totient
kernel, or the totient's prime-dilation laws.

\subsection{An unlabelled algebraic branch condition loses arithmetic information}
There are infinitely differing sequences $g:\N_{>0}\to\N$ with
$g(n)-\ph(n)=o(n)$, $g(n)=\ph(n)$ for odd $n$, and
\[
 (g(2n)-g(n))(g(2n)-2g(n))=0\qquad(n\ge1).
\]
To construct one, let $t_j$ be the product of the first $j$ odd primes.
Put $g(2t_j)=2\ph(t_j)$ and $g(n)=\ph(n)$ elsewhere.
The odd parts $t_j$ are distinct, so the altered dyadic chains do not overlap.
At $n=t_j$, the multiplier is $2$; at $n=2t_j$, it is $1$; thereafter the
usual multiplier $2$ is restored. Every other chain is unchanged.
Moreover
\[
 \frac{g(2t_j)-\ph(2t_j)}{2t_j}
 =\frac12\prod_{\substack{p\mid t_j}}(1-1/p)\longrightarrow0.
\]
This proves the assertions and $0\le g(n)\le n$.
No rationality assertion is made about this new control's series.

For comparison, the \emph{labelled} laws
\[
 g(pn)=\begin{cases}p\,g(n),&p\mid n,\\(p-1)g(n),&p\nmid n\end{cases}
\]
for even one fixed prime $p$, together with $g-\ph=o(n)$, force $g=\ph$
on positive integers. If a defect at $n$ were nonzero, iteration along
$p^j n$ would keep its ratio to the argument nonzero. This elementary
rigidity is already Lean-checked in
\texttt{PrefixValuationAndControlRigidity.lean}; it is not a new result here.
The example above shows why forgetting which local class chooses each branch
weakens the information.

\subsection{Bounded modulus does not pass automatically to the parent}
For $A_m=\sum_{n\ge1}(\ph(n)\bmod m)2^{-n}$, the elementary bound
$\ph(n)\le n-1$ for $n\ge2$ gives
\[
 0<S-A_m\le\sum_{n=m+1}^{\infty}(n-1)2^{-n}
 =(m+1)2^{-m}.
\]
Thus the bounded projections approach $S$ very rapidly, but irrationality or
joint linear independence of all those projections gives no separation from
a fixed rational candidate at this shrinking scale.
The certificate \eqref{eq:lower-linear} is controlled by the centre $N$ of
an auxiliary pulse, not by the original modulus $m$.

The already formalised termwise-window obstruction makes the lost uniformity
concrete. If $N+L>1$ and $2^L\mid\ph(N+L)$, then
\[
 2^L\le\ph(N+L)\le N+L-1,
 \qquad \frac{N+L+2}{2^L}>1.
\]
Consequently an envelope-only argument that asks this tail bound to be less
than $1/q$ cannot work for any integer $q\ge1$.
This statement concerns that explicit termwise divisibility scheme; it does
not exclude cancellations in a weighted totient prefix.

\section{The remaining arithmetic endpoint and formal interfaces}
The live canonical equivalence is retained. For every $c\ge0$ and positive odd
$v$, one must find $H>0$ with $\ph(v)\mid H$ and, writing
\[
 M=(2^H-1)/v,\qquad
 B_{H,c}=\sum_{j=0}^{H-1}\ph(c+1+j)2^{H-1-j},
\]
one must prove
\[
 c+H+1<(-B_{H,c})\bmod M<M-(c+H+1).
\]
No universal producer for this condition is obtained here. The new kernel
algebra and value-classification theorems are subsidiary results with complete
ordinary proofs, not substitutes for the quantified residue gap.

\paragraph{Formalisation plan.}
The accompanying candidate\par
{\small\ttfamily ProgressionLocalTotientIndependence.lean}\par
extends the compiled periodic transport lemma to coefficients that freeze on
some progression through each point, and to eventual relations.
{\small\ttfamily FiniteDilationPulseCore.lean} supplies the minimum-divisor and
neighbour-divisibility arithmetic for the new pulse construction.
{\small\ttfamily TotientDyadicPolynomialIdentity.lean} gives the quadratic relation from
the two exact dyadic laws, without claiming its completeness.
None was compiled in this environment. The density lemma, polynomial ideal
completeness, series rearrangement and the full CRT prime producer still require
formalisation. The existing isolated-pulse theorem is the appropriate binary
consumer; the ordinary proof here also covers every positive integer base.

\paragraph{Checks.}
The checker uses exact totients, modular row reduction with integer arithmetic,
finite CRT witnesses and exact rational bounds. Its separate receipt records
what was tested. It makes no numerical inference about irrationality or about
an unbounded rank from a finite matrix. The existing Lean build, Comparator
replay and large continued-fraction computation were not rerun.

\paragraph{Status and antecedents.}
The source authority is the supplied round-six packet, whose live short note has
title \emph{A Basis for the 2-Kernel of Euler's Totient} and principal result
\texttt{thm:kkernelrank}. That theorem is not replaced here.
The density construction belongs to the Schinzel tradition. Garcia--Udell--Yu
\cite{guy} describe the classical multidimensional results and impose additional
prime-tuple conditions in their own conditional theorems. Halupczok--Ohst
\cite{ho} establish affine one-coordinate and quotient density statements.
Martin \cite{martin} gives a stronger antecedent for the affine \emph{linear}
independence in the existing note. Coons \cite{coons} already proves non-regularity.
No priority is claimed for the density argument or its algebraic consequence.
The exact principal-ideal normal form and the joint value classification below
were not located in the sources checked; this is not a priority determination.
No new theorem in this memorandum has been compiled in Lean.

\begin{thebibliography}{9}\small
\bibitem{martin} G. Martin, \emph{Simultaneous inequalities among values of the
Euler phi-function}, 2006, \href{https://arxiv.org/abs/math/0603053}{arXiv:math/0603053}.
\bibitem{coons} M. Coons, \emph{(Non)Automaticity of number theoretic functions},
J. Th\'eor. Nombres Bordeaux \textbf{22} (2010), 339--352,
\href{https://doi.org/10.5802/jtnb.718}{doi:10.5802/jtnb.718}.
\bibitem{guy} S. R. Garcia, G. Udell and J. Yu,
\emph{Multidimensional Schinzel-type theorems for multiplicative functions near
prime tuples}, J. Number Theory \textbf{219} (2021), 212--227,
\href{https://arxiv.org/abs/2004.13208}{arXiv:2004.13208}.
\bibitem{ho} K. Halupczok and M. Ohst, \emph{Density properties of fractions with
Euler's totient function}, Involve \textbf{19} (2026), 1--31,
\href{https://arxiv.org/abs/2411.11065}{arXiv:2411.11065}.
\bibitem{wong} E. Wong, answer to \emph{An infinite sum based on the mod-parity
of Euler's totient function}, Mathematics Stack Exchange, 29 March 2015,
\href{https://math.stackexchange.com/questions/1210886}{question 1210886}.
Wong's denominator base equals the modulus; the joint theorem above fixes a
common underlying base and mixes its distinct powers.
\bibitem{erdos} P. Erd\H{o}s, \emph{On arithmetical properties of Lambert series},
J. Indian Math. Soc. (N.S.) \textbf{12} (1948), 63--66.
The signed bounded-pulse formulation in the supplied corpus isolates the
separation mechanism used here.
\bibitem{packet} Supplied round-six \#249 packet, \texttt{01\_short\_note.tex},
\texttt{03\_research\_packet.json}, and the modules
\texttt{ResidueClassTotientSeries.lean},
\texttt{AllBaseTotientKernel.lean},
\texttt{ParityPerturbedRationalControl.lean},
\texttt{PrefixValuationAndControlRigidity.lean} and
\texttt{PeriodicTotientIndependence.lean}.
\end{thebibliography}
\end{document}
